% ============================================================
% CHƯƠNG 6 — KẾT LUẬN VÀ HƯỚNG PHÁT TRIỂN  (mẫu SoICT chương 6)
% Đổi số chương 5 -> 6 ngày 17/11/2026 (tệp cũ Chap5.tex) khi Chương 2–3 đổi thành Phân tích, Thiết kế.
% Tạo lại ngày 07/11/2026 theo mẫu Do_An_Latex_SoiCT. Chương 5 cũ "KẾT QUẢ THỰC NGHIỆM" đã gộp vào Chương 3
% (bản cũ: backup/2026-11-07_truoc-mau-soict/Chap5.tex); khung Kết luận trước đây nằm trong main.tex.
%
% Hướng dẫn của mẫu SoICT:
%   - Kết luận: so sánh kết quả với các sản phẩm, nghiên cứu tương tự; đã làm được gì, chưa làm được gì,
%     đóng góp nổi bật, bài học kinh nghiệm;
%   - Hướng phát triển: trước hết là công việc để hoàn thiện chức năng đã làm, sau đó là hướng mới để
%     cải thiện, nâng cấp.
% Văn phong: xem ghi chú đầu Chap1.tex (hạn chế xưng "em"; mục con dùng \subsection*, không \ref tới mục con).
% ============================================================

\chapter{KẾT LUẬN VÀ HƯỚNG PHÁT TRIỂN}
\label{chap:ketluan}

% [VIẾT SAU] Các mục dưới đây chưa viết nội dung — bỏ dấu % để khôi phục.
% \section{Kết luận}
% \label{sec:ketluan}
% % Đối chiếu từng mục tiêu ở mục \ref{sec:muctieu}: đạt / chưa đạt (smart contract, backend, giao diện,
% % kiểm thử >= 80%, Sepolia). Bảng so sánh với Verra Registry, UNDP, CAD Trust, Toucan theo tiêu chí: quyền sở hữu
% % on-chain, chống tính toán kép, mô hình dữ liệu chuẩn, mã nguồn mở, kiểm tra công khai.
% % Hạn chế đã biết từ mã nguồn: chưa xử lý reorg khi đồng bộ sự kiện; chưa kiểm thử tải API; NFT gửi vào pool
% % bằng transferFrom (không phải safeTransferFrom) sẽ bị kẹt; backend lưu tệp cục bộ (chưa IPFS); dữ liệu phát thải
% % tự khai báo; tên trường CAD Trust cần đối chiếu từ điển dữ liệu V2.0 chính thức.
% % Bài học kinh nghiệm.
% % [VIẾT SAU] \canviet{Kết luận}
%
% \section{Hướng phát triển}
% \label{sec:huongphattrien}
% % Hoàn thiện: sửa các hạn chế trên, kiểm toán bảo mật, IPFS, xử lý reorg.
% % Hướng mới: lớp MRV số hóa với IoT (Vaccargiu \cite{vaccargiu2026blockchain}); chữ ký số của VVB; ánh xạ sang
% % mã trong nước và số sê-ri theo Thông tư 11/2026/TT-BNNMT; liên thông CAD Trust để chặn phát hành kép giữa các
% % hệ thống; xuất báo cáo AEF cho Điều 6.2; bảo vệ quyền riêng tư (Liu \cite{liu2024blockchain}).
% % [VIẾT SAU] \canviet{Hướng phát triển}
